\section*{Capítulo \ref{ch:b1:cross-sections} --- Seções geológicas}

\begin{solution}{exo:b1:cross-sections:1}
Na escala 1:25\,000, \qty{1}{cm} no mapa correspondem a \qty{250}{m} no terreno. A escala vertical é de \qty{1}{cm} para \qty{100}{m}, logo a escala vertical é
$250/100 = 2.5$ vezes a horizontal: $e = 2.5$.
\end{solution}

\begin{solution}{exo:b1:cross-sections:2}
$\gamma = \ang{90} - \ang{60} = \ang{30}$, e $\tan\beta = \tan\ang{40}
\cos\ang{30} = 0.8391\times0.8660 = 0.7267$: $\beta = \ang{36.0}$, a camada
descendo em direção à extremidade leste da seção. Na seção de \omterm{def:b1:spherical-astronomy:horizontal}{azimute}
\ang{000}, $\gamma = \ang{90}$ e $\beta = 0$: a seção corre ao longo da \omterm{def:b1:geological-maps:attitude}{direção} e a camada aparece horizontal.
\end{solution}

\begin{solution}{exo:b1:cross-sections:3}
$t = h\cos\delta = 120\times\cos\ang{12.16} = 120\times0.9776 =
\qty{117.3}{m}$.
\end{solution}

\begin{solution}{exo:b1:cross-sections:4}
$\tan\ang{12} = 0.2126$. Com $e = 2$: $\arctan 0.4251 = \ang{23.0}$; com
$e = 3$: $\arctan 0.6377 = \ang{32.5}$; com $e = 5$: $\arctan 1.0628 =
\ang{46.7}$. O ângulo cresce bem menos que o próprio $e$, e nunca chega a
\ang{90}.
\end{solution}

\begin{solution}{exo:b1:cross-sections:5}
(a) $t = 300\sin\ang{25} = \qty{126.8}{m}$. (b) $t = 300\sin\ang{35} =
\qty{172.1}{m}$. (c) $t = 300\sin\ang{15} = \qty{77.6}{m}$. A mesma largura de \omterm{def:b1:geological-maps:geological-map}{afloramento} pode significar espessuras que diferem em mais de duas vezes: a inclinação do terreno deve sempre ser levada em conta.
\end{solution}

\begin{solution}{exo:b1:cross-sections:6}
Com \omterm{def:b1:geological-maps:attitude}{sentido de mergulho} $\alpha$ e \omterm{def:b1:geological-maps:attitude}{mergulho} $\delta$, o teorema dá, na seção norte--sul ($\theta = 0$), $\tan\ang{20} = \tan\delta\cos\alpha$,
e na seção leste--oeste ($\theta = \ang{90}$) $\tan\ang{30} =
\tan\delta\cos(\alpha - \ang{90}) = \tan\delta\sin\alpha$. Dividindo,
$\tan\alpha = 0.5774/0.3640 = 1.586$, $\alpha = \ang{057.8}$; somando os
quadrados, $\tan\delta = \sqrt{0.3640^2 + 0.5774^2} = 0.6825$, $\delta =
\ang{34.3}$. Os dois \omterm{def:b1:cross-sections:apparent-dip}{mergulhos aparentes} são menores que o \omterm{def:b1:geological-maps:attitude}{mergulho} verdadeiro, como devem ser.
\end{solution}

\begin{solution}{exo:b1:cross-sections:7}
$100/\tan\ang{12.16} = 100/0.2155 = \qty{464.1}{m}$. No mapa, as \omterm{def:b1:cross-sections:structure-contour}{curvas estruturais} do flanco leste estão de fato um pouco menos de meio quilômetro afastadas, igualmente espaçadas e paralelas: o flanco é um plano.
\end{solution}

\begin{solution}{exo:b1:cross-sections:8}
B1: $524.02 - 488.98 = \qty{35.04}{m}$; B3: $529.98 - 494.93 =
\qty{35.05}{m}$. Acima de B2, a base do folhelho é o topo do calcário, \qty{180}{m} acima da sua base: $280.16 + 180 = \qty{460.16}{m}$,
mais alto que o terreno (\qty{380.64}{m}), por isso ele foi erodido ali.
Entre B1 e B2 o contato sobe $460.16 - 35.04 = \qty{425.12}{m}$
em \qty{2000}{m}: uma inclinação de $0.2126$, um \omterm{def:b1:cross-sections:apparent-dip}{mergulho aparente} de \ang{12.0}. B1
e B3, à mesma distância de cada lado de B2, encontram todos os \omterm{def:b1:geological-maps:geological-map}{contatos} na mesma altitude com uma precisão de \qty{1}{cm}: a \omterm{def:b1:geological-maps:fold}{dobra} é simétrica em torno de um plano axial vertical em $x = \qty{3.5}{km}$, com os dois flancos mergulhando igualmente, e nenhuma falha corta a seção entre eles.
\end{solution}

\begin{solution}{exo:b1:cross-sections:9}
\omterm{def:b1:cross-sections:throw}{Afastamento horizontal} $= 150/\tan\ang{50} = 150/1.1918 = \qty{125.9}{m}$. Na seção de
\omterm{def:b1:spherical-astronomy:horizontal}{azimute} \ang{150}, $\gamma = \ang{30}$ e $\tan\beta = 1.1918\times0.8660 =
1.0321$: a falha aparece com o \omterm{def:b1:cross-sections:apparent-dip}{mergulho aparente} \ang{45.9}. Os cortes de uma camada horizontal estão separados por \qty{150}{m} na vertical sobre esse traço, logo
$150/1.0321 = \qty{145.3}{m}$ de separação ao longo da seção, mais que o \omterm{def:b1:cross-sections:throw}{afastamento horizontal}, já que a seção cruza a falha obliquamente.
\end{solution}

\begin{solution}{exo:b1:cross-sections:10}
$\cos\gamma = \tan\ang{20}/\tan\ang{40} = 0.3640/0.8391 = 0.4338$, logo
$\gamma = \ang{64.3}$: as seções de \omterm{def:b1:spherical-astronomy:horizontal}{azimutes} $\alpha \pm \ang{64.3}$
(e as direções opostas, onde a camada sobe). O \omterm{def:b1:cross-sections:apparent-dip}{mergulho aparente} é
\ang{40.0}, \ang{39.6}, \ang{38.3} e \ang{36.0} para $\gamma = 0$, \ang{10},
\ang{20}, \ang{30}, mas \ang{22.8}, \ang{16.0}, \ang{8.3} e $0$ para
$\gamma = \ang{60}$, \ang{70}, \ang{80}, \ang{90}. De fato, derivando
$\tan\beta = \tan\delta\cos\gamma$,
$\dfrac{\mathrm d\beta}{\mathrm d\gamma} =
-\dfrac{\tan\delta\sin\gamma}{1 + \tan^2\delta\cos^2\gamma}$, que é nula
em $\gamma = 0$ e vale $-\tan\delta$ em $\gamma = \ang{90}$: uma seção um pouco afastada da \omterm{def:b1:geological-maps:attitude}{direção} do \omterm{def:b1:geological-maps:attitude}{mergulho} mostra quase o \omterm{def:b1:geological-maps:attitude}{mergulho} verdadeiro, enquanto perto da \omterm{def:b1:geological-maps:attitude}{direção}, cada grau de \omterm{def:b1:spherical-astronomy:horizontal}{azimute} muda o \omterm{def:b1:cross-sections:apparent-dip}{mergulho aparente} em quase um grau.
\end{solution}

\begin{solution}{exo:b1:cross-sections:11}
As camadas mergulham para leste, logo a normal ao acamamento, seguida para baixo, está inclinada \ang{12.16} da vertical para oeste, em direção a \ang{260.7}.
(a) Um furo vertical faz com ela o ângulo $\psi = \ang{12.16}$:
$176.0/\cos\ang{12.16} = \qty{180.0}{m}$, a \omterm{def:b1:cross-sections:thickness}{espessura vertical}. (b) Um
furo que cai \ang{60} em direção a \ang{260.7} está inclinado \ang{30} da vertical do mesmo lado da normal: $\psi = \ang{30} - \ang{12.16} =
\ang{17.84}$, $\cos\psi = 0.9519$, e o testemunho tem $176.0/0.9519 =
\qty{184.9}{m}$ de comprimento. (c) Em \omterm{def:b1:geological-maps:attitude}{direção} a \ang{080.7} ele está inclinado do outro lado:
$\psi = \ang{30} + \ang{12.16} = \ang{42.16}$, $\cos\psi = 0.7413$, e o
testemunho tem $176.0/0.7413 = \qty{237.4}{m}$ de comprimento. O testemunho é o mais curto, igual a $t$, quando o furo corre ao longo da normal, caindo \ang{77.84} em direção a \ang{260.7}.
\end{solution}

\begin{solution}{exo:b1:cross-sections:12}
No plano vertical leste--oeste, medindo os ângulos a partir do leste, as camadas superiores sobem para leste com \ang{20} e as inferiores descem para leste com \ang{70}. Girando o \omterm{def:b1:geological-maps:geological-map}{afloramento} em torno da \omterm{def:b1:geological-maps:attitude}{direção} de \ang{20}, com o lado leste para baixo, as camadas superiores ficam horizontais e as inferiores passam a $\ang{70} + \ang{20} = \ang{90}$: estavam verticais quando as camadas superiores foram depositadas. História, a partir do evento mais antigo: as camadas inferiores são depositadas na horizontal; são inclinadas até a vertical e soerguidas; suas bordas são erodidas até formar uma superfície; as camadas superiores são depositadas sobre essa superfície; o conjunto é inclinado \ang{20} para oeste; a erosão corta a superfície do terreno de hoje. É a história do \omterm{def:b1:geological-maps:geological-map}{afloramento} de Hutton em Siccar Point.
\end{solution}

\begin{solution}{pb:b1:cross-sections:1}
\textbf{1.} A espessura acima do túnel é $z - 200$: a maior em
$x = \qty{1}{km}$, $556 - 200 = \qty{356}{m}$; a menor em $x =
\qty{9}{km}$, $339 - 200 = \qty{139}{m}$.

\textbf{2.} \qty{1}{km} de altitude mede \qty{4.2}{cm} e \qty{1}{km}
de distância mede \qty{1.4}{cm}: $e = 4.2/1.4 = 3$.

\textbf{3.} $380.64 - 220.48 = \qty{160.16}{m}$.

\textbf{4.} Em B2 a base do calcário está em $380.64 - 100.48 =
\qty{280.16}{m}$ e a base do arenito em \qty{160.16}{m}: o
túnel, a \qty{200}{m}, está no arenito.

\textbf{5.} B1: $524.02 - 668.98 = \qty{-144.96}{m}$; B3: $529.98 - 674.93 =
\qty{-144.95}{m}$.

\textbf{6.} De oeste para leste: giz, folhelho, calcário, arenito (no núcleo
do \omterm{def:b1:geological-maps:fold}{anticlinal}), calcário, folhelho e de novo giz até
$x = \qty{9}{km}$, passando pelo \omterm{def:b1:geological-maps:fold}{sinclinal} e atravessando a falha.

\textbf{7.} $\gamma = \ang{80.7} - \ang{90} = \ang{-9.3}$; $\tan\beta =
\tan\ang{12.16}\cos\ang{9.3} = 0.2155\times0.9869 = 0.2126$: $\beta =
\ang{12.0}$, descendo para leste.

\textbf{8.} $(280.16 + 144.95)/2000 = 425.11/2000 = 0.2126$: \ang{12.0}.

\textbf{9.} $\arctan(3\times0.2126) = \arctan 0.6377 = \ang{32.5}$.

\textbf{10.} $\gamma = \ang{120} - \ang{90} = \ang{30}$; $\tan\beta =
\tan\ang{60}\cos\ang{30} = 1.732\times0.866 = 1.500$: $\beta = \ang{56.3}$.

\textbf{11.} Os flancos mergulham para \ang{080.7} ou \ang{279.3}, apenas
\ang{9.3} da linha da seção, e $\cos\ang{9.3} = 0.987$; a falha mergulha para
\ang{120}, \ang{30} dela, e $\cos\ang{30} = 0.866$.

\textbf{12.} O plano da falha é $z = 250 - \tan\ang{60}\,[(x - 4300)
\sin\ang{120} + (y - 1000)\cos\ang{120}]$. Com $y = 6000$, $(y - 1000)
\cos\ang{120} = -2500$, logo $(x - 4300)\sin\ang{120} = 2500 - (z - 250)/
\tan\ang{60}$, e $1/\tan\ang{60} = \tan\ang{30}$.

\textbf{13.} Em $z = 200$: $x = 4300 + (2500 + 50\times0.5774)/0.8660 =
4300 + 2528.9/0.8660 = \qty{7220}{m}$, isto é, \qty{7.22}{km} de C.

\textbf{14.} \omterm{def:b1:cross-sections:throw}{Afastamento horizontal} $= 200/\tan\ang{60} = \qty{115.5}{m}$; separação ao longo de
C--D $= 200/1.500 = \qty{133.3}{m}$.

\textbf{15.} A falha mergulha para leste-sudeste e é normal: o teto, a leste, está abatido. No nível do túnel o giz está dos dois lados da falha: a oeste dela, o giz do \omterm{def:b1:geological-maps:fold}{sinclinal}; a leste, giz rebaixado em \qty{200}{m}.

\textbf{16.} O túnel permanece no giz através da falha, de modo que nenhuma mudança de rocha na frente a anuncia; apenas uma zona de rocha triturada, e talvez água, o faz. Os engenheiros preveem a sua posição a partir da seção (aqui \qty{7.22}{km}), perfuram sondagens de avanço à frente da máquina conforme ela se aproxima e preparam o suporte para a zona triturada.

\textbf{17.} $668.98 - 488.98 = \qty{180.00}{m}$.

\textbf{18.} $t = 180\cos\ang{12.16} = \qty{176.0}{m}$.

\textbf{19.} Na crista ($x = 3500$) a base do arenito está em
\qty{160.16}{m} (questão 3); ela desce dos dois flancos com o \omterm{def:b1:cross-sections:apparent-dip}{mergulho aparente}
\ang{12.0} da questão 7, isto é, de $\tan\ang{12.0}$ por metro de
distância horizontal a partir do eixo.

\textbf{20.} A base do calcário está \qty{120}{m} acima da base do
arenito, em $280.16 - u\tan\ang{12.0}$ com $u = |x - 3500|$, e o seu topo,
\qty{180}{m} mais alto, em $460.16 - u\tan\ang{12.0}$. Com $\tan\ang{12.0} =
0.21256$: a base está abaixo de \qty{200}{m} quando $u > 80.16/0.21256 = 377.1$,
e o topo acima dela quando $u < 260.16/0.21256 = 1224.0$.

\textbf{21.} $x = 3500 \pm 377.1$ e $3500 \pm 1224.0$: o túnel está no calcário
de \qty{2276}{m} a \qty{3123}{m} e de \qty{3877}{m} a \qty{4724}{m}.

\textbf{22.} Ao longo de uma linha horizontal, a base e o topo do calcário, \qty{180}{m} afastados na vertical e ambos inclinados de $\tan\beta$, são cruzados a uma distância horizontal de $h/\tan\beta$: $180/0.21256 = \qty{846.8}{m}$ em cada
flanco, o mesmo, a menos do arredondamento, que $1224.0 - 377.1 = 846.9$.

\textbf{23.} $2\times846.8 = \qty{1693.7}{m}$: o túnel perfura
\textbf{cerca de \qty[text-series-to-math=true]{1694}{m} de calcário}, mais de um quinto de seus
\qty{8}{km} de comprimento, em dois trechos de \qty{847}{m} de cada lado do núcleo de arenito do \omterm{def:b1:geological-maps:fold}{anticlinal}.
\end{solution}
